\documentclass[journal]{IEEEtran}
% --- encoding & fonts (pdflatex) ---
\usepackage[utf8]{inputenc}
\usepackage[T1]{fontenc}
\usepackage{textcomp}
% --- tables/figures (pandoc body needs these) ---
\usepackage{graphicx}
\usepackage{array}
\usepackage{booktabs}
\usepackage{longtable}
\usepackage{calc}
\providecommand{\real}[1]{#1}
\providecommand{\tightlist}{\setlength{\itemsep}{0pt}\setlength{\parskip}{0pt}}
\usepackage{url}
\usepackage{cite}
\hyphenation{op-tical net-works semi-conduc-tor Industry}

\begin{document}

\title{Digital Twins in Engineering and STEM Higher Education: A Systematic Review of Architectures, Pedagogy, and the Quality of Learning Evidence}

\author{Segundo~F.~Segura-Altamirano,~Diana~M.~Castro-Cárdenas,~Carlos~L.~Oblitas-Vera,~Grimaldo~D.~Benavides-Campos,~Janet~R.~Aquino-Lalupú,~Wilder~A.~Alvarado-Castillo,~Rosa~E.~Salazar-Cabrejos,~and~Nilton~C.~German~Reyes%
\thanks{All authors are with the Universidad Nacional Pedro Ruiz Gallo, Lambayeque, Peru (corresponding author: Segundo F. Segura-Altamirano; e-mail: sseguraal@unprg.edu.pe).}}

\markboth{IEEE Transactions on Learning Technologies, Manuscript submitted 2026}%
{Segura-Altamirano \MakeLowercase{\textit{et al.}}: Digital Twins in Engineering and STEM Higher Education}

\maketitle

\begin{abstract}
Practical laboratory training is central to engineering education yet constrained by cost, safety, access, and post-pandemic continuity, and digital twins (DTs)---virtual replicas bidirectionally coupled to physical assets---are increasingly proposed to address these limits. This article systematically reviews how DTs are designed, implemented, and evaluated in higher-education engineering and STEM settings, across four objectives: technological architecture and deployment environments; pedagogy and Industry 4.0/5.0 competencies; empirical learning outcomes; and institutional benefits and barriers. Following PRISMA 2020, five databases were searched; a terminological screen, three-agent large-language-model semantic screening with a moderator (inter-agent Fleiss' $\kappa = 0.60$), and human review yielded 126 eligible records, of which 65 retrievable studies (2018--2026) formed the corpus. The methodological quality of empirical studies was scored on a transparent 0--6 composite and sensitivity-tested, and the synthesis was read through an instructional-alignment lens (constructive alignment, ICAP, and Kirkpatrick's levels). The corpus demonstrates feasibility, engagement, and disciplinary breadth, with a recent shift toward bidirectional, occasionally AI-augmented twins; however, the technical maturity of the twin is decoupled from the rigour of its evaluation---architectures are described more than instrumented, pedagogy is active but seldom named, and the learning evidence is weak-dominated, with only a minority of rigorous studies reporting substantial effects. We map where transfer-oriented, well-reported evidence is most needed and propose a reporting checklist for educational digital twins.
\end{abstract}

\begin{IEEEkeywords}
Constructive alignment, digital twin, engineering education, evidence quality, Industry 4.0, laboratory learning, learning outcomes, STEM education.
\end{IEEEkeywords}

\section{Introduction}\label{introduction}

Laboratory and hands-on instruction occupy a privileged place in engineering
education. The instructional laboratory is not an adjunct to the lecture but a
setting with its own learning objectives---instrumentation, experimentation,
modelling, design, and the disciplined translation of theory into practice---that
the field has articulated explicitly and treated as constitutive of professional
competence \cite{feisel_role_2005,lei_teaching_2022,naukkarinen_supporting_2018}.
Yet the physical laboratory is a constrained resource. Time, space, equipment
cost, maintenance, and safety place hard limits on how many students can practise,
how often, and on how current the available apparatus remains relative to the
industrial state of the art \cite{lei_teaching_2022,prichetnikov_investigation_2020}.
The COVID-19 pandemic sharpened these long-standing constraints, forcing the rapid
migration of practical courses online and exposing how poorly procedural skill
formation tolerates the loss of physical access \cite{haase_teaching_2022,alkhedher_interactive_2021}. In response, virtual, remote and simulation-based
laboratories have proliferated across engineering disciplines---from electrical
measurement, transmission lines and microcontroller programming to chemical
reaction engineering and control systems \cite{chumchuen_remote_2023,andrieu_remote_2021,panagiotakis_remote_2022,naukkarinen_supporting_2018,lei_teaching_2022}.

The displacement of practical work from physical apparatus to its digital
surrogates raises a question that has organised engineering-education research for
two decades: do non-traditional laboratories achieve the same learning outcomes as
hands-on work? Comparative reviews have repeatedly found that virtual and remote
laboratories can match hands-on settings on many measured outcomes, with the
relative advantage depending on the outcome and the instructional design rather
than on the medium alone \cite{ma_handson_2006,corter_constructing_2007,lindsay_effects_2005,brinson_learning_2015}. That literature is equally
notable, however, for what it says about evidence: Brinson's review of the
empirical research reports that the great majority of studies find ``no significant
difference'' yet are dominated by weak designs and self-report, and bibliometric
syntheses document rapid growth in publication without commensurate growth in
methodological rigour \cite{brinson_learning_2015,heradio_virtual_2016,potkonjak_virtual_2016}. A recent review of digital technologies for practical
teaching reaches the same double conclusion: these tools improve access,
flexibility and engagement, but the evidence that they replace the tacit,
procedural learning of the physical workshop remains limited
\cite{maladzhi_digital_2025}. The digital twin enters engineering education as the
most recent and most technically ambitious member of this lineage, and it inherits
the lineage's unresolved question.

Among these digital approaches, the digital twin (DT) has emerged as a distinct
and more demanding proposition. First articulated by Grieves in the
context of product-lifecycle management as reported in \cite{liljaniemi_using_2020}, a DT is a virtual representation of a physical asset
maintained in correspondence with it through a data connection, such that the
state of one is reflected in the other across the asset's lifecycle
\cite{de_larmelina_conceptualizing_2026,weichbroth_toward_2024}. This definitional core is widely shared, but the term
remains contested. Analysing more than 15,000 publications, \cite{abdelrahman_what_2025} show that no consensus definition exists and that several features often
assumed to be intrinsic---real-time bidirectional data flow, simulation, and
embedded artificial intelligence (AI)---are not yet mature even in industrial
practice. The literature therefore distinguishes degrees of integration, from a
static digital model, through a digital shadow with one-way physical-to-virtual
updating, to a fully bidirectional digital twin capable of closed-loop control
Kritzinger et al., 2018, as cited in \cite{liljaniemi_using_2020,abdelrahman_what_2025,bokhtiar_al_zami_digital_2025}. This gradation matters pedagogically, because it
separates a DT from a conventional simulation or an immersive scene: what is
labelled a ``twin'' may range from a high-fidelity but disconnected model to an
instrumented, synchronised cyber-physical system \cite{kabir_digital_2025,lee_digital_2025}. Establishing this maturity---rather than assuming it from the
label---is consequently a prerequisite for any serious account of what educational
DTs are and what they can teach.

The motivation to bring DTs into the classroom is inseparable from the competence
demands of Industry 4.0 and the emerging, human-centred Industry 5.0
\cite{saini_industry_2026,kabir_digital_2025}. Employers and accreditation bodies
increasingly expect graduates to be fluent in the industrial Internet of Things,
automation and data analytics and, with the Industry 5.0 turn, in human--machine
collaboration and sustainability-oriented practice \cite{sell_international_2025,azofeifa_top_2025,dautaj_importance_2026}. Several studies document a
persistent mismatch between graduate competencies and labour-market expectations
\cite{denissova_factors_2026}, while survey evidence shows that engineering faculty
themselves often feel insufficiently prepared, and insufficiently equipped, to
teach these technologies \cite{cuperman_academic_2026}. DTs are proposed in this
context precisely because they can expose students to the same
cyber-physical systems they will meet in industry, and align that exposure
with competence frameworks articulated for Industry 5.0 work and learning
factories \cite{zare_engineering_2026,nitu_engineering_2026,al-jumeily_implementing_2026}.

DTs rarely arrive alone. They are typically embedded within, or compared against,
a broader wave of immersive and virtual technologies---virtual, augmented and mixed
reality, the metaverse, and serious games---whose use in engineering and higher
education has expanded rapidly and been mapped by numerous reviews
\cite{lampropoulos_virtual_2025,jin_extended_2026,fan_impacts_2025,espinoza_virtual_2025,llanos-ruiz_virtual_2025}. The evidence on these
technologies is, however, mixed. Many studies report gains in motivation,
engagement and the understanding of complex concepts, but controlled comparisons
sometimes find no significant difference in measured learning outcomes relative
to traditional methods, underscoring that pedagogical design, not technology
alone, drives the effect \cite{giussani_study_2025,yindeemak_emerging_2026,lainidis_systematic_2026}. Within this landscape, DT-based education has
grown from early demonstrations in construction and mechanical engineering
\cite{sepasgozar_digital_2020,liljaniemi_using_2020,treffinger_investigations_2022}
into a varied corpus spanning manufacturing, robotics, electrical and chemical
engineering, aviation maintenance and even university mathematics
\cite{khan_teaching_2026,kwateng_enhancing_2026,hunpinyo_isa95_2026,xiang_research_2025,bolanos_design_2026,lee_digital_2023}, often coupled with
project-based and experiential pedagogies and framed as preparation for
Industry 4.0/5.0 practice \cite{zhang_effectiveness_2024,ruppert_demonstration_2023,terkaj_virtual_2024,alvarez_ariza_investigating_2026}.

This momentum has produced several reviews, but they leave the lineage's
central question---whether the technology achieves the learning outcomes of
hands-on work, and on what quality of evidence---largely unanswered for DTs.
Existing syntheses tend either to survey immersive technologies broadly,
with DTs as one instance among many \cite{lampropoulos_virtual_2025,tokhayeva_virtualisation_2026,jiang_state_2026}, or to concentrate on a single
application niche, such as DT-coupled virtual learning environments for smart
manufacturing \cite{kumar_karanam_systematic_2025}. Across this body of work three
gaps recur. First, the techno-pedagogical architecture of educational DTs---their
deployment environment, communication protocols and synchronisation, and degree
of bidirectional integration---is seldom characterised jointly with the pedagogy it
is meant to serve, so that technical and instructional accounts run in parallel
rather than being read together. Second, and continuing the evidence-quality
concern that comparative laboratory reviews have raised since Brinson
\cite{brinson_learning_2015,heradio_virtual_2016}, reviews rarely appraise the
methodological quality of the underlying learning evidence; effectiveness claims
are aggregated without weighting how rigorously each was established, even though
the immersive-technology literature itself warns that headline gains often rest
on small, uncontrolled designs \cite{fan_impacts_2025,giussani_study_2025}.
Third, no synthesis has examined how technical, pedagogical and evaluative
choices co-vary across the field---whether, for example, the most technically
capable twins are also the most rigorously evaluated, or whether ambitious
competence claims are matched by demonstrated outcomes. Addressing these gaps
calls for a synthesis that treats architecture, pedagogy, evidence quality and
institutional context as linked dimensions rather than as separate narratives,
and that reads them through an explicit instructional lens rather than as a
technical inventory.

Accordingly, this systematic review synthesises 65 primary studies on the use of
digital twins in engineering and higher education, published between 2018 and
2026. The corpus was coded through an inductive thematic analysis organised into
eight themes, and the synthesis is then read through the alignment lens set out in
the next section. On that basis the review pursues four objectives, each posed as
a question about the instructional system rather than the artefact: (O1) to
characterise the technological architectures and deployment environments of
educational DTs---their enabling components, synchronisation and convergence with
AI---and the instructional affordances these provide; (O2) to describe the
pedagogical designs adopted and the Industry 4.0/5.0 competencies they target, and
to ask whether that pedagogy engages students constructively and aligns
competencies to the twin; (O3) to appraise the empirical evidence for learning
outcomes---at what level of evaluation, and with what methodological rigour, it
rests; and (O4) to assess the institutional benefits claimed for DTs and the
barriers to their adoption, and whether those claims are aligned with demonstrated
outcomes. Beyond profiling each dimension separately, the review analyses how
these dimensions co-occur and relate, so that the maturity of the technology can
be read against the rigour of its evaluation and the breadth of competence claims
against the weight of demonstrated learning. By coupling a structured description
of what is being built with a critical appraisal of what has actually been shown,
the review aims to give engineering educators, researchers and institutions an
evidence-weighted map of the field and a foundation for the research agenda
developed in the sections that follow.

The significance of this work for engineering education research and practice lies
in its appraisal stance. Rather than cataloguing digital-twin applications, the
review weighs the \emph{quality} of the evidence behind them and asks how technical,
pedagogical and evaluative choices co-vary, reading the field as an instructional
system through the alignment lens developed in the next section rather than as a
catalogue of technology. The intent is to give researchers an evidence-weighted
basis for identifying where confirmatory, transfer-oriented work is most needed,
and to give educators and institutions a more cautious basis for adopting digital
twins---pairing them with explicit pedagogy, aligned assessment, and instrumented
cost, sustainability and accessibility claims---than a technology-led reading would
support. What the synthesis shows about the balance between technical maturity and
evaluative rigour, and what follows from it, is developed in the Results and
Discussion.

\section{Related Reviews}\label{related-reviews}

Several reviews already touch this field, which makes it important to state
precisely what they do and do not establish. Table 1 compares the most relevant
prior syntheses---those addressing digital twins (DTs) in education and the longer
line of comparative laboratory reviews from which the DT question descends---along
the dimension that motivates the present work: whether each weighs the
\emph{methodological quality} of the learning evidence it aggregates, rather than only
cataloguing or counting studies.

%\clearpage
%\onecolumn
\footnotesize
\begin{table}[!t]
	\renewcommand{\arraystretch}{1.2}
	\setlength{\tabcolsep}{3pt}
	\caption{Prior Reviews of Digital Twins and Technology-Mediated Laboratories}
	\label{tab:T1}
	\centering
	\scriptsize
	\begin{tabular}{@{} l p{3.0cm} c p{3.2cm} @{}}
		\hline
		\textbf{Review} & \textbf{Focus and Scope} & \textbf{Qual.?} & \textbf{Stated Gap / Limitation} \\ \hline
		\cite{lampropoulos_virtual_2025} & VR in engineering ed. (342 studies) & No & No methodological quality analysis \\
		\cite{tokhayeva_virtualisation_2026} & Virtualisation in HE; tools, cost & Partial & Weak primary designs; limited inference \\
		\cite{jiang_state_2026} & VR in HE (1,100 papers) & No & Little empirical work on learning effects \\
		\cite{kumar_karanam_systematic_2025} & DT virtual env. for manufacturing & No & No cost or scalability appraisal \\
		\cite{chamorro-atalaya_use_2024} & DT in university teaching (92 papers) & No & No analysis of governance / competence \\
		\cite{omrany_digital_2026} & DT for construction training (19 studies) & No & Fragmented, largely conceptual literature \\
		\cite{vergara_decade_2025} & DT in materials sci. \& eng. (1,106 papers) & No & Educational outcomes not appraised \\
		\cite{maladzhi_digital_2025} & Digital tech. for practical SET teaching & Yes & Tools do not fully replace physical labs \\
		\cite{ma_handson_2006} & Hands-on vs. simulated vs. remote labs & No & Taxonomy framework; no quality weighting \\
		\cite{brinson_learning_2015} & Learning outcomes: non-traditional labs & Partial & Weak empirical research designs dominate \\
		\cite{heradio_virtual_2016} & Virtual and remote labs mapping & No & Primary-study quality not appraised \\
		\cite{potkonjak_virtual_2016} & Virtual labs for science and engineering & No & Reviewed without formal risk-of-bias \\ \hline 
		\multicolumn{4}{@{}l}{\parbox{\linewidth}{\vspace{1mm}\textit{Note:} HE = higher education; SET = science, engineering and technology; DT = digital twin; Qual.? = Appraises learning evidence quality.}}
	\end{tabular}
\end{table}

%\setlength{\tabcolsep}{4pt}
%\begin{longtable}{@{}
%  >{\raggedright\arraybackslash}p{(\linewidth - 8\tabcolsep) * \real{0.1700}}
%  >{\raggedright\arraybackslash}p{(\linewidth - 8\tabcolsep) * \real{0.1200}}
%  >{\raggedright\arraybackslash}p{(\linewidth - 8\tabcolsep) * \real{0.3300}}
%  >{\raggedright\arraybackslash}p{(\linewidth - 8\tabcolsep) * \real{0.1300}}
%  >{\raggedright\arraybackslash}p{(\linewidth - 8\tabcolsep) * \real{0.2500}}@{}}
%\caption{Prior reviews of digital twins and of technology-mediated laboratories in higher education, and whether each appraised the methodological quality of the underlying learning evidence. HE = higher education; SET = science, engineering and technology; DT = digital twin.}\label{tab:T1}\\
%\toprule\noalign{}
%\begin{minipage}[b]{\linewidth}\raggedright
%Review
%\end{minipage} & \begin{minipage}[b]{\linewidth}\raggedright
%Type
%\end{minipage} & \begin{minipage}[b]{\linewidth}\raggedright
%Focus and scope
%\end{minipage} & \begin{minipage}[b]{\linewidth}\raggedright
%Appraises evidence quality?
%\end{minipage} & \begin{minipage}[b]{\linewidth}\raggedright
%Stated gap
%\end{minipage} \\
%\midrule\noalign{}
%\endhead
%\bottomrule\noalign{}
%\endlastfoot
%\cite{lampropoulos_virtual_2025} & Scoping & VR in engineering education (342 studies) & No & No methodological-quality analysis of included studies \\
%\cite{tokhayeva_virtualisation_2026} & Systematic & Virtualisation in HE; tools and cost & Partial & Weak primary designs; causal and economic inference limited \\
%\cite{jiang_state_2026} & Bibliometric & VR in HE (1,100 papers) & No & Little empirical work on educational effects \\
%\cite{kumar_karanam_systematic_2025} & Systematic & DT-coupled virtual environments for manufacturing education & No & No cost or scalability appraisal \\
%\cite{chamorro-atalaya_use_2024} & Bibliometric & DT in university teaching and learning (92 papers) & No & No analysis of governance or teacher competence \\
%\cite{omrany_digital_2026} & Systematic & DT for construction education and training (19 studies) & No & Fragmented, largely conceptual literature \\
%\cite{vergara_decade_2025} & Bibliometric & DT in materials science and engineering (1,106 papers) & No & Output mapped; educational outcomes not appraised \\
%\cite{maladzhi_digital_2025} & Systematic (PRISMA) & Digital technologies for practical SET teaching (28 studies) & Yes & Digital tools do not fully replace physical laboratories \\
%\cite{ma_handson_2006} & Comparative review & Hands-on vs simulated vs remote laboratories & No & Establishes taxonomy and comparison; no quality weighting \\
%\cite{brinson_learning_2015} & Review of empirical research & Learning outcomes: non-traditional vs hands-on labs & Partial & Prevailing ``no significant difference'', but weak designs dominate \\
%\cite{heradio_virtual_2016} & Bibliometric & Virtual and remote labs & No & Trends mapped; primary-study quality not appraised \\
%\cite{potkonjak_virtual_2016} & Review & Virtual laboratories for science, technology and engineering & No & Effectiveness reviewed without formal risk-of-bias assessment \\
%\end{longtable}
\normalsize



Two patterns are visible. First, the reviews that \emph{do} weigh evidence quality
(e.g., the PRISMA review of digital practical-teaching technologies, and the
nursing and mobile-VR meta-analyses that apply GRADE-style appraisal
\cite{maladzhi_digital_2025,li_effects_2026,widowati_meta-analysis_2025}) sit
\emph{outside} the digital-twin literature. Second, every review \emph{specific to digital
twins in education}---whether bibliometric \cite{chamorro-atalaya_use_2024,vergara_decade_2025} or systematic \cite{kumar_karanam_systematic_2025,omrany_digital_2026}---characterises or counts the literature without appraising the
rigour of the learning evidence it reports. To our knowledge, then, the present
review is the first synthesis of digital twins in engineering and STEM higher
education to make the \emph{quality} of the learning evidence its organising axis, and to
read architecture, pedagogy and evidence jointly rather than as separate
inventories---directly extending, for digital twins, the evidence-quality concern
that comparative laboratory reviews have raised since Ma \& Nickerson and Brinson
\cite{ma_handson_2006,brinson_learning_2015,heradio_virtual_2016}.

\section{Conceptual Framework}\label{conceptual-framework}

The objection that motivates this section is methodological as much as
substantive. A catalogue of digital-twin (DT) implementations, however complete,
does not by itself tell an engineering educator whether these systems teach, or
how. To read architecture, pedagogy and evidence as parts of a single
instructional system rather than as separate technical and educational
narratives, we adopt three established lenses from the learning sciences and use
them consistently to interpret the synthesis. They are deliberately modest in
number and chosen because each maps onto one face of the review's question.

\textbf{Constructive alignment} \cite{biggs_enhancing_1996,biggs_teaching_2011} supplies
the organising principle. In Biggs' formulation, a course functions as a system
in which intended learning outcomes, the teaching and learning activities, and
the assessment tasks are mutually aligned: students learn what the activities and
assessment actually require of them, not what the syllabus declares. A DT is, in
these terms, a learning \emph{activity}---often an elaborate one. The pertinent question
is therefore not how sophisticated the twin is, but whether the intended
competencies and the assessment have been aligned with it. This reframing turns
the review's central tension---technical maturity advancing faster than evaluative
rigour---into a question about systemic alignment rather than about technology, and
it connects the corpus to a long-standing concern in engineering education with
the explicit, assessable purposes of laboratory work \cite{feisel_role_2005}.

\textbf{The ICAP framework} \cite{chi_icap_2014} refines how we read the pedagogy of
Objective O2. Decades of evidence establish that active engagement outperforms
passive reception in engineering and STEM instruction \cite{prince_does_2004,freeman_active_2014}, yet ``active learning'' is used so broadly that it obscures
more than it reveals. ICAP distinguishes four modes of cognitive engagement of
rising expected benefit---passive, active, constructive and interactive---on the
basis of what learners observably do with material. Applied to the corpus, it
lets us ask whether DT activities labelled ``active'' engage students
\emph{constructively} (generating new artefacts or explanations) and \emph{interactively}
(co-constructing through dialogue), or merely keep them behaviourally busy. The
single most capable engagement mode in ICAP---learners constructing the model
itself---is precisely the activity our corpus documents only once, a gap the
framework makes legible.

\textbf{Kirkpatrick's levels of evaluation} \cite{kirkpatrick_kirkpatricks_2016} structure
the appraisal of learning outcomes in Objective O3. The four levels---reaction,
learning, behaviour and results---order evidence by how consequential and how
difficult to obtain it is. Reactions (satisfaction, perceived usefulness) are
the easiest to collect and the least probative; demonstrated behaviour change and
downstream results are the most demanding. This ordering provides a vocabulary
for the review's finding that affective and reaction-level outcomes dominate the
corpus while transfer and retention---the levels that would substantiate the
field's competence claims---are rarely measured.

Together these lenses let the review pose, for educational DTs, the comparative
question that has organised research on non-traditional laboratories for two
decades: whether technology-mediated practical work achieves the learning
outcomes of hands-on work, and on what quality of evidence \cite{ma_handson_2006,brinson_learning_2015}. Accordingly, the four objectives are read as alignment
questions. O1 asks what instructional affordances the architecture actually
provides; O2 asks whether the pedagogy engages students constructively and aligns
competencies to the twin; O3 asks at what Kirkpatrick level, and with what rigour,
learning has been evidenced; and O4 asks whether institutional claims are aligned
with demonstrated outcomes or run ahead of them.

A note on the status of this framework is essential to its honest use. The corpus
was coded inductively: the eight themes (T1--T8) that structure the extraction
were derived from the studies themselves, not deduced from constructive
alignment, ICAP or Kirkpatrick. These three lenses were adopted at the
\emph{synthesis} stage, to interpret the extracted data and to connect the findings to
engineering-education research; they were not imposed on the protocol or used as
an a priori coding scheme. The framework therefore guides interpretation and
discussion, and the mapping of pedagogical labels to ICAP modes and of outcomes
to Kirkpatrick levels is offered as a reading of the evidence, not as a
measurement claimed to be present in the primary studies.

\section{Positionality Statement}\label{positionality-statement}

The authors are engineering and education faculty at a public university, with
teaching and research interests spanning electronic and control engineering,
automation, and engineering education. Several of us have designed and taught
laboratory courses using simulation- and technology-mediated tools, an
experience that shaped this review in two ways. First, it motivated our central
question --- whether the technological sophistication of a digital twin is matched
by evidence that it improves learning --- because we have seen compelling
demonstrations adopted ahead of evidence in our own teaching contexts. Second, it
disposed us toward a deliberately cautious, evidence-weighted reading rather than
an advocacy stance. We hold no commercial interest in any digital-twin platform,
vendor, or engine. We are mindful that our location in a resource-constrained
setting sharpens our attention to cost, accessibility, and equity, and may make
us more sceptical of unquantified institutional-benefit claims than authors
writing from better-resourced contexts. To limit the influence of these
dispositions on the findings, study selection used independent screening agents
with explicit inclusion/exclusion tests, and the quality appraisal followed a
pre-specified rubric applied with double verification.

\section{Methods}\label{methods}

\subsection{Protocol and reporting}\label{protocol-and-reporting}

This review was designed and is reported following the PRISMA 2020 statement and,
for the conduct of a systematic literature review in engineering education
specifically, the methodological guidance synthesised by \cite{borrego_systematic_2014}.
A pre-specified protocol (Review Protocol v2) fixed the research question,
eligibility criteria, search strategy, screening procedure and analysis plan
before screening began. The review was not prospectively registered; registration
in PROSPERO or OSF is recommended before any update. The protocol, extraction
schemas and analysis code are openly archived, so the design is auditable in full. The review question, framed with the
Population--Concept--Context (PCC) scheme, asked how digital twins (DTs) are
designed, implemented and evaluated in higher-education engineering and STEM
settings. It was decomposed into four objectives: technological architectures and
deployment environments (O1); pedagogical strategies and Industry 4.0/5.0
competencies (O2); empirical evidence of learning outcomes (O3); and institutional
benefits and critical barriers (O4).

\subsection{Eligibility criteria}\label{eligibility-criteria}

The population and context were university-level engineering or STEM students and
instructors in courses, laboratories, training or assessment; the concept was the
educational implementation, deployment, evaluation or use of a DT. The construct
was defined canonically and applied with a deliberately conservative severity rule
to guard against semantic drift: a DT is a virtual representation of a \emph{specific}
physical asset, bidirectionally and continuously coupled to it, so that state
changes propagate physical\ensuremath{\rightarrow}virtual \emph{and} virtual\ensuremath{\rightarrow}physical. Standalone simulations
or numerical models, virtual or remote laboratories, VR/AR/XR, the metaverse,
serious games, generic CPS/IoT testbeds, digital models (no automatic data
exchange) and digital shadows (one-way coupling) were treated as near-misses and
qualified only when the bidirectional coupling was explicit. Studies that claimed
a ``digital twin'' with partial or unstated synchronisation were coded \emph{uncertain}
rather than excluded, under an authoritative precedence rule whereby absence of
confirming evidence is never grounds for exclusion.

A hard document-type gate was applied before thematic relevance: empirical journal
articles, full conference papers and early-access empirical articles were
eligible, whereas reviews of any kind, surveys, editorials, opinion pieces,
abstracts, posters and keynote summaries were ineligible. No language restriction
was imposed; the retained corpus includes English, Spanish and Russian.

\subsection{Information sources and search}\label{information-sources-and-search}

Five bibliographic databases were searched: Scopus, Web of Science, IEEE Xplore,
EBSCO and Dimensions. The searches were last run on 14 May 2026 (Web of
Science, IEEE Xplore) and 15 May 2026 (Scopus, EBSCO, Dimensions). The core string combined the construct with an education
block --- (``digital twin*'' OR ``educational digital twin*'') AND (``engineering
education'' OR ``engineering training'' OR ``technical education'' OR ``STEM
education'') --- adapted to each database's syntax. The choice of a compact,
precision-oriented string rests on a documented sensitivity analysis of the
search chain. In Scopus, the construct alone (``digital twin*'') returned 20,917
records, and lengthening the construct disjunction with ``digital-twin*'' or
``educational digital twin*'' did not enlarge this domain; conjoining the
education block reduced the set to 192--193 records, whereas appending a long
technology disjunction (laboratory, IoT, CPS, mechatronics, automation,
telecommunications, embedded and cyber-physical systems) only narrowed recall
further, to 159, without surfacing additional relevant work. Progressively
trimming the education terms moved the yield monotonically between 141
(``engineering education'' alone) and 193 (the full education disjunction). These
runs confirmed two design decisions: the term ``digital twin*'' alone retrieves
the same domain as longer construct disjunctions, and the education block --- not a
long technology disjunction --- is the effective constraint. The string was
therefore fixed to favour precision over an unfocused expansion of recall.

The searches returned 1,504 raw records (Scopus 772, Dimensions 409, IEEE 181,
EBSCO 81, Web of Science 61). After removing 38 records lacking a title or
authors, 1,466 records entered deduplication. Matching on DOI (232) and fuzzy
title (47) removed 279 duplicates, yielding 1,187 unique records; a further 9
without an abstract were set aside, leaving 1,178 records for screening.

\subsection{Study selection}\label{study-selection}

Selection proceeded in three stages, summarised in the PRISMA 2020 flow diagram
(Figure 1). First, an automated terminological screen required both a DT block
and an engineering-education block to be present in the title, abstract or
keywords, with a keyword-based rescue for records carrying the terms only in
author or index keywords. This retained 799 records (206 by full field match and
593 by keyword rescue) and set aside 379 on terminological grounds. Second, the
799 records underwent semantic screening by a panel of three independent large
language model (LLM) reviewers. To avoid one-dimensional judgments, the three
reviewers shared a single base model (DeepSeek-V4-Flash, a non-reasoning
configuration at temperature 0) but were instantiated with three complementary
expert \emph{roles}: a Principal Investigator (judging the genuine intersection of all
required concepts), a Research Methodologist (document type, empirical design and
analytic purpose), and a Domain Expert (construct centrality and population/context
authenticity). Each role applied the \emph{same} pre-specified, decidable criterion
set---the canonical construct definition with its strict / moderate / near-miss
severity mapping, the hard document-type gate, the positive include/exclude tests,
and an authoritative uncertainty-precedence rule (absence of confirming evidence is
coded \emph{uncertain}, never \emph{exclude})---but weighted it through its professional lens;
the role defines the lens, not the scope. Each reviewer returned, for every record,
an independent decision (include, exclude or uncertain), a confidence score, the
governing criterion, and a verbatim evidence span. Running the base model at
temperature 0 with cached calls makes the procedure \emph{reproducible}---re-running
yields identical outputs---but it does not make the three reviewers identical: their
differing role lenses produce genuine, criterion-grounded disagreement on
borderline records, which is the intended behaviour of a review panel rather than
stochastic noise. Aggregation followed a strict-unanimity rule: a record was
auto-classified only when all three reviewers agreed; any split, or any ``uncertain''
vote, was escalated to a fourth moderator agent (same model, independent prompt and
seed) that read the three rationales and, absent a defensible resolution, deferred
the record to human review.

Agreement among the three primary reviewers was substantial: they were unanimous on
466 of the 799 records (58\%), the moderator resolved a further 96 split cases (5 to
include, 91 to exclude), and inter-agent reliability was Fleiss' \ensuremath{\kappa} = 0.60 across the
three categories, with pairwise Cohen's \ensuremath{\kappa} between 0.56 and 0.63. This \ensuremath{\kappa} quantifies
\emph{agreement among the role lenses, not the validity of the screen}: it is expected to
fall below unity precisely because the panel is designed to disagree on genuinely
borderline records, and because the uncertainty-precedence rule routes ambiguous
records to the ``uncertain'' category---depressing nominal agreement---rather than risking
a false exclusion. The decision criteria were anchored to twelve expert-authored
calibration cases (four include, four exclude, four uncertain) spanning the
construct-versus-near-miss boundary. Validity against a human standard was secured at
the inclusion margin---every record the panel did not unanimously exclude (247) was
adjudicated by a human---while residual false-negative risk in the unanimously-excluded
zone is bounded by a documented, reproducible audit (a blind random sample of
consensus-excluded records re-adjudicated by a human from title, abstract and
keywords, with the models' rationales withheld to prevent anchoring); the audit
script is released in the repository to be completed before any update. The agent
prompts, scores and per-record outputs are provided as supplementary material and in
the public repository. Third, a human
review of the 247 records that the agents flagged as include (11) or uncertain
(236) confirmed 126 as eligible. Full texts could be
obtained for 65 of these; the remaining 61 could not be retrieved, and the
synthesis proceeds on the 65 retrievable studies.

\begin{figure}[!t]
\centering
\includegraphics[width=\columnwidth,keepaspectratio]{figs/F1_prisma.png}
\caption{PRISMA 2020 flow diagram of identification, screening and inclusion. Title/abstract screening combined an automated terminological screen with semantic screening by three independent large language model (LLM) agents and a fourth moderator agent (inter-agent Fleiss' \ensuremath{\kappa} = 0.60).}
\label{fig:F1}
\end{figure}

\subsection{Data extraction}\label{data-extraction}

Full texts were converted from PDF with the Marker pipeline to preserve textual
fidelity, and extraction followed a theme-conditional activation principle: each
study activated only the extraction templates of the themes in which it had been
inductively coded, with the remaining variables left as not-applicable. This
reduced extraction load, kept the evidence pertinent, and allowed per-theme
reporting rates to be computed. Every critical variable used a controlled
vocabulary --- NR (not reported), NA (not applicable) and UNCERTAIN --- and double
verification, recording the extracted value together with a verbatim supporting
quotation and its location, so that each datum is traceable to its source and
independently auditable.

\subsection{Thematic analysis and synthesis}\label{thematic-analysis-and-synthesis}

An inductive thematic analysis --- conducted with three frontier LLMs (ChatGPT 5.5,
Gemini 3.1 Pro and Claude Opus 4.7), whose independent codings were iteratively
cross-fed and refined through successive rounds of author review to converge on a
scheme granular enough to be exhaustive yet parsimonious enough to remain
tractable --- produced eight themes: practical infrastructure (T1), enabling
architectures (T2), AI convergence (T3), Industry 4.0/5.0 competencies (T4),
pedagogical design (T5), empirical evidence (T6), institutional value (T7) and
barriers (T8). The three models' outputs are included in the public repository. These were mapped N:M onto the four objectives
through a traceability matrix (O1 = T1+T2+T3; O2 = T4+T5; O3 = T6; O4 = T7+T8), so
that a single study could contribute to several themes and objectives.

The synthesis was structured and narrative, without meta-analysis --- a decision
justified by the methodological and outcome heterogeneity of the corpus. Because
the review's principal inferences depend on the strength of the primary evidence,
the methodological quality of each empirical study was formally appraised rather
than assumed. Because no single instrument captures the rigour markers that the
engineering-education research community has long argued distinguish trustworthy
from anecdotal evidence \cite{streveler_conducting_2006}, we used a transparent,
purpose-built composite rather than importing a clinical appraisal tool wholesale:
each empirical study scored 0--6 on five additive markers---design rigour,
comparator, inferential statistics, effect size and declared validity
threats---chosen because they are the dimensions most consistently flagged in that
methodological literature and are extractable from the reporting actually present
in the corpus. The composite is deliberately simple and its banding is tested for
robustness rather than treated as exact (below). Analysis proceeded in three
principal layers: per-theme descriptive profiling with explicit
reporting rates (Layer 1); aggregation by objective, including the 0--6
evidence-quality composite over empirical studies
(Layer 2); and theme co-occurrence and configuration analysis (Layer 3). These
were complemented by an exploratory relational layer comprising an
attribute-association network (\ensuremath{\phi} coefficients with Fisher exact tests and
non-parametric bootstrap 95\% confidence intervals, B = 5,000), a multiple
correspondence analysis embedding, and a capability-by-evidence gap map. The
robustness of the evidence bands was tested with a threshold-sensitivity analysis
and bootstrap intervals, and the full pipeline is scripted end-to-end for
reproducibility.

\section{Results}\label{results}

This section reports the synthesis of 65 primary studies on digital twins (DTs)
in engineering education. Following the review design, the analysis proceeds in
three layers --- descriptive profiling per theme (Layer 1), aggregation by review
objective (Layer 2), and cross-cutting co-occurrence and relational analysis
(Layer 3/4) --- mapped onto the four objectives O1--O4. Counts are reported over
the subset of studies that activated each theme; reporting rates (the share of
those studies that actually documented a given variable, as opposed to ``not
reported'', NR) are stated explicitly because they bound the strength of every
claim that follows. All supporting figures (Figures 2--11) and the
synthesis tables (Tables 2 and 3) are produced by a reproducible analysis pipeline.

\subsection{Corpus overview}\label{corpus-overview}

The corpus spans 2018--2026 and is recent: 39 of 65 studies (60\%) were published
in 2025--2026, with 26 in 2025 and 13 in 2026. Journal articles (37) and
conference papers (26) are represented in comparable numbers. Studies are
unevenly distributed across the eight inductive themes: pedagogy (T5, n=47) and
Industry 4.0/5.0 competencies (T4, n=47) are most frequent, followed by
practical infrastructure (T1, n=38), empirical evidence (T6, n=35), barriers
(T8, n=28), enabling architectures (T2, n=26), institutional value (T7, n=21)
and AI convergence (T3, n=11). Most studies activated three (19) or four (28)
themes, indicating that individual papers typically address technology,
pedagogy and outcomes jointly rather than in isolation.

\subsection{O1 --- Technological architectures and deployment environments (T1+T2+T3)}\label{o1-technological-architectures-and-deployment-environments-t1t2t3}

\textbf{Deployment environment (T1).} The environment type was reported in 37 of 38
studies (97\%). Hybrid physical--virtual setups are the most common configuration
(18), followed by immersive VR/AR/MR environments (11), remote (4) and
virtual-only (3). The modelled physical system was almost always specified
(95\%), most often a manufacturing plant or production line, with isolated cases
in metallurgy, instrumentation laboratories, HVAC and robotics. Two variables
were comparatively under-reported: the access modality (37\%) and the immersive
hardware/software stack (55\%), where Unity is the recurring engine and head-set
hardware (e.g., Oculus Rift) appears only sporadically. Reported visual/
functional fidelity (58\%) was qualitatively ``high'' in 14 studies, but fidelity
was seldom quantified.

\textbf{Enabling architecture (T2).} Among the 26 architecture studies, physical
components, virtual components and communication protocols were each documented
in 88\% of cases, whereas the layered architecture was described in 62\%. OPC UA
is the most frequently named protocol, with MQTT, TCP/IP, EtherCAT, the Asset
Administration Shell (AAS) and the IEC 61131-3/61499 family also present. Mapped
onto the Kritzinger et al.~taxonomy (reported by 18 of 26), the studies that
specify a category claim a fully bidirectional digital twin in 10 cases and a
(unidirectional) digital shadow in 7, with a single static digital model.
Synchronisation, where stated (20 of 26), is ``real-time'' (8) or ``near-real-time''
(8); only one study documents automatic bidirectional control. These figures
should be read with caution: the taxonomy is largely self-described by authors
and rarely accompanied by latency or jitter measurements, so the prevalence of
``complete'' twins likely reflects terminological usage more than verified
closed-loop synchronisation.

\textbf{AI convergence (T3).} AI is the least represented dimension (11 studies).
Where an AI family is named (8 of 11), it ranges across classical machine
learning, deep/neural networks, generative models (LLM/GAN), explainable AI and
knowledge graphs; the most common declared functions are prediction, diagnosis
and predictive maintenance, with a few cases of automated student feedback and
adaptive personalisation. Notably, AI performance metrics (accuracy, F1,
inference latency) were reported in only 3 of 11 studies (27\%), so the
educational and technical performance of these AI components is largely
unverified in the current corpus.

\begin{figure}[!t]
\centering
\includegraphics[width=\columnwidth,keepaspectratio]{figs/F1_O1_arquitectura.png}
\caption{Technological architectures and deployment environments (Objective O1): category distributions and per-variable reporting rates for deployment environment, modelled system, enabling architecture and AI convergence. AR = augmented reality; MR = mixed reality; VR = virtual reality; AAS = Asset Administration Shell; AI = artificial intelligence; OPC UA = Open Platform Communications Unified Architecture.}
\label{fig:F2}
\end{figure}

\subsection{O2 --- Pedagogical design and competencies (T4+T5)}\label{o2-pedagogical-design-and-competencies-t4t5}

\textbf{Disciplinary scope and competencies (T4).} The engineering domain was
reported in 46 of 47 studies (98\%). Manufacturing (12) and automation (9)
dominate, but the corpus is disciplinarily diverse, also covering mechatronics
(5), chemical engineering (4), and smaller clusters in metallurgy, energy,
civil/construction, robotics and cybersecurity. Targeted competencies were
documented in 81\% of studies and cluster around IIoT, smart manufacturing,
automation, data analytics and --- increasingly --- human--machine collaboration
framed as Industry 5.0. Two alignment variables are weak spots: explicit mapping
to industrial or accreditation standards (ABET, ISA-95, RAMI 4.0, IEC) and
documented industry--academia collaboration were each reported in only 21\% of
studies. Competencies are therefore frequently \emph{asserted as design intent}
rather than anchored to external frameworks or industry validation.

\textbf{Instructional design (T5).} The pedagogical model was reported in 43 of 47
studies (92\%) and is overwhelmingly described as ``active learning'' (23) or
``experiential learning'' (14); formally named approaches such as project-based
learning, learning/teaching factories, ADDIE or challenge-based learning are
individually rare. Interventions are concentrated at the undergraduate level
(31), and student participation is typically ``active experimentation'' (25) or a
``full life-cycle project'' (12), with student \emph{construction} of the twin itself
documented only once. The teacher is consistently positioned as a facilitator or
supervisor (reported in 77\%). A recurring reporting gap is intervention
duration, stated in only 38\% of studies, which limits any inference about dose
or sustained effects.

Read through the ICAP lens \cite{chi_icap_2014}, the reported participation modes
cluster at the \emph{active} level of cognitive engagement: the dominant mode, ``active
experimentation'' (25 of 42), denotes hands-on manipulation of the twin rather than
the generation of new artefacts or explanations. Modes that would reach the
\emph{constructive} and \emph{interactive} levels---where ICAP predicts the largest learning
gains---are the minority: full life-cycle projects (12) and, at the apex, students
\emph{constructing} the twin itself, which appears once. Guided exploration (4) sits at
the active/passive boundary. This distribution is consistent with a field that has
adopted the language of active learning \cite{prince_does_2004,freeman_active_2014}
without yet designing for its more demanding, generative forms; the mapping is an
interpretive reading of the extracted labels rather than a measure reported by the
primary studies.

\begin{figure}[!t]
\centering
\includegraphics[width=\columnwidth,keepaspectratio]{figs/F2_O2_pedagogia.png}
\caption{Pedagogical design and Industry 4.0/5.0 competencies (Objective O2): engineering domains, targeted competencies, pedagogical models and participation modes, with reporting rates. IIoT = industrial Internet of Things.}
\label{fig:F3}
\end{figure}

\subsection{O3 --- Empirical evidence of learning outcomes (T6)}\label{o3-empirical-evidence-of-learning-outcomes-t6}

Thirty-five studies activated T6. Research design was reported in 80\% of them
and is mixed in rigour: quasi-experiments (8) and conceptual/design-development
studies (8) are equally common, alongside case studies/surveys/pilots (8),
mixed-methods (3) and a single controlled experiment. Comparator use is best read
in three categories rather than as a single ``reported'' rate: of the 35 studies,
20 (57\%) employed an actual comparator, 9 (26\%) explicitly had none, and 6 (17\%)
did not report it. Among the 20 with a comparator, the most common were a
traditional class (10) or a within-subject pre/post design (7), with a few using
a no-treatment control group. Critically for a technology whose promise is to
stand in for hands-on practice, only one study compared the twin against a
physical laboratory \cite{liljaniemi_enhancing_2025}---the very comparison the field's
foundational question requires \cite{ma_handson_2006,brinson_learning_2015}. Across
reported outcome dimensions, affective/engagement results
(23; spanning motivation, technology acceptance, self-efficacy and, where
NASA-TLX is used, cognitive load) and cognitive results (22) are reported more often than procedural skills
(17).

The central limitation of the evidence base concerns statistical reporting.
Inferential statistics with reported values appear in only 11 of 35 studies
(31\%), and an effect size in only 8 (23\%). To make this explicit, we scored
each empirical study on a 0--6 composite (design rigour, comparator, inferential
statistics, effect size, and declared validity threats). The resulting distribution is weighted toward
the lower end: 20 studies are \emph{weak}, 10 \emph{moderate} and 5 \emph{strong} (median
composite score = 2 of 6). This bottom-heavy profile does not depend on the
chosen cut-points: under a stricter banding only 2 studies are strong and under a
lenient banding 12 are, yet weak-or-moderate studies remain the clear majority
in every scheme (66--94\%; 20 of 35 score \ensuremath{\leq}2 and are weak under all schemes).

Where strong studies exist, they report substantial effects. In a controlled
quasi-experiment on a digital-twin/IoT laboratory intervention (N=200: 140
intervention, 60 control), \cite{kwateng_enhancing_2026} report a between-groups difference of t(198)=8.94, p\textless.001, a
pre/post mean gain of +18.5 points in the intervention group versus +4.3 in the
control, a large effect on sensor-integration skills (Cohen's d=0.91), and
procedural improvements such as cloud-synchronisation success
rising from 68\% to 92\%; \cite{wang_innovative_2025} report Cohen's d=1.24 for knowledge
mastery, \cite{tarng_application_2024} an ANCOVA group effect of \ensuremath{\eta^{2}}=0.50, and \cite{lu_digital_2025} Cohen's d=1.06 in algorithm-design and d=0.91 in debugging tasks. These cases show
that rigorous evaluation of educational DTs is feasible, and the five strongest
studies are not small (N=60--200). They remain the exception, however: they are
mostly single-institution, non-randomised quasi-experiments with short, single-module
exposures and possible novelty effects --- limitations the authors themselves
declare in 57\% of T6 studies. Reported gains should therefore be read as
promising but not yet robustly generalisable.

\begin{figure}[!t]
\centering
\includegraphics[width=\columnwidth,keepaspectratio]{figs/F3_O3_evidencia.png}
\caption{Empirical evidence of learning outcomes (Objective O3): research designs, comparators, reported outcome dimensions and the 0--6 evidence-quality composite. NASA-TLX = NASA Task Load Index.}
\label{fig:F4}
\end{figure}

\subsection{O4 --- Institutional benefits and critical barriers (T7+T8)}\label{o4-institutional-benefits-and-critical-barriers-t7t8}

\textbf{Institutional value (T7).} Among the 21 studies addressing institutional
value, accessibility and equity arguments are the most frequently articulated
(86\%), predominantly framed as geographic remote access; the geographic context
of application was reported in all 21. However, the quantitative
underpinnings of the institutional case are thin: implementation cost was
reported in only 6 studies (29\%), environmental-sustainability indicators in 3
(14\%), and licensing model in 8 (38\%, of which 6 open-source). The few studies
that do quantify these dimensions are informative: \cite{bonavolonta_sustainability_2026}
provide a life-cycle assessment of an immersive electronics laboratory
(\ensuremath{\approx}7,000 kWh/year, \ensuremath{\approx}1.75 tCO\textsubscript{2}/year, \ensuremath{\approx}200 kg of WEEE over seven years, and a
per-student cost roughly four times lower than the traditional laboratory, with
completion rising to 70\% vs 56\%), and \cite{gallego_low-cost_2025} report a low-cost
VR-teleoperated robotic arm at under €400. Post-COVID continuity, a common
narrative motivation, is explicitly invoked in only 38\% of these studies.

\textbf{Barriers (T8).} Twenty-eight studies discuss barriers, but reporting is again
partial. Technical barriers are the most frequently documented (43\%) ---
insufficient model fidelity, limited interoperability, latency/jitter,
synchronisation complexity and cybersecurity --- while institutional/teacher
barriers (teacher training, curricular governance, resistance to change) and
pedagogical/curricular barriers (the virtual-to-physical transfer gap, limited
ecological validity, misalignment of competencies--activities--assessment) are
each documented in roughly a quarter of studies. Scalability or transferability
evidence was reported in 36\%, usually as architectural argument rather than
demonstrated multi-institution replication. Declared facilitating conditions
(25\%) converge on institutional commitment, industry partnerships, prior
training and minimum infrastructure.

\begin{figure}[!t]
\centering
\includegraphics[width=\columnwidth,keepaspectratio]{figs/F4_O4_beneficios_barreras.png}
\caption{Institutional benefits and barriers (Objective O4): articulated benefits (accessibility, cost, sustainability, continuity) and reported barrier types (technical, institutional, pedagogical), with reporting rates. WEEE = waste electrical and electronic equipment.}
\label{fig:F5}
\end{figure}

\subsection{Cross-cutting analysis: co-occurrence, gaps and relationships}\label{cross-cutting-analysis-co-occurrence-gaps-and-relationships}

The co-occurrence matrix (Figure 6) shows that the corpus is organised around a
pedagogy--competencies--evidence core: the strongest theme pairings are T4--T5
(competencies \ensuremath{\times} pedagogy, 34 studies) and T5--T6 (pedagogy \ensuremath{\times} evidence, 31). AI
(T3) is weakly connected to every other theme (maximum co-occurrence of 8, with
T4). This structure motivates four configuration analyses that surface
relationships not visible in the per-theme frequencies:

\begin{itemize}
\tightlist
\item
  \textbf{Aspiration--evidence gap.} Twenty-five studies (38\%) claim Industry 4.0/5.0
  competencies (T4) without activating empirical evaluation (T6), and 16 (25\%)
  describe a pedagogical design (T5) with no empirical evidence. The competency
  claims of the field currently outpace its demonstrated learning outcomes.
\item
  \textbf{Optimism asymmetry.} Sixteen studies (25\%) report institutional benefits
  (T7) without discussing any barrier (T8), while 23 report barriers without
  claiming benefits. The presence of a substantial barriers-only group suggests
  the corpus is not uniformly promotional, but the benefits-without-barriers
  group should be read with caution.
\item
  \textbf{Immersive-without-evidence.} Eight studies (12\%) combine an immersive or
  practical environment (T1) with a pedagogical design (T5) but report no
  empirical outcomes (T6) --- a recurring pattern of compelling demonstrations
  awaiting evaluation.
\item
  \textbf{Integrated maturity is rare.} Only 3 studies (5\%) assemble the full stack
  of environment + architecture + pedagogy + evidence (T1+T2+T5+T6:
  \cite{sakkas_multiplayer_2025}; \cite{tarng_application_2024}; \cite{terkaj_framework_2024}), and
  only 1 combines a detailed architecture, AI and empirical evaluation
  (T2+T3+T6: \cite{machado_automatic_2025}).
\end{itemize}

The relational tables make three further associations explicit; given the small
cell sizes these are descriptive, not inferential.

\begin{enumerate}
\def\labelenumi{\arabic{enumi}.}
\tightlist
\item
  \textbf{Architecture maturity is decoupled from evaluation rigour.} Only 4 studies
  both characterise the twin on the Kritzinger taxonomy (T2) \emph{and} report an
  evaluable outcome (T6) with a quality band. In other words, the studies that
  describe \emph{how complete} the twin is are largely not the studies that test
  \emph{whether it improves learning} --- a structural separation between the
  engineering and the educational evidence in the field.
\item
  \textbf{Effect-size reporting is low across all designs.} The share of studies
  reporting an effect size is 0\% for conceptual/design-development and
  mixed-methods studies, about 38\% for case-study/survey/pilot and
  quasi-experimental studies, and 100\% only for the single controlled
  experiment (Figure 7a). Even quasi-experiments report an effect size in only
  about a third of cases, so weak statistical reporting is pervasive rather than
  confined to the weakest designs.
\item
  \textbf{A shift toward fully bidirectional twins over time.} Among architecture
  studies, ``complete'' digital twins are concentrated in 2024--2026 (10 of 18),
  whereas earlier years are dominated by digital shadows (Figure 7b). The
  terminological and technical ambition of the field is increasing, which makes
  the decoupling noted in point 1 --- and the scarce reporting of AI and
  synchronisation metrics --- the more consequential to address.
\end{enumerate}

\begin{figure}[!t]
\centering
\includegraphics[width=\columnwidth,keepaspectratio]{figs/F5_coocurrencia_temas.png}
\caption{Theme co-occurrence matrix: number of studies jointly activating each pair of the eight inductive themes (T1--T8).}
\label{fig:F6}
\end{figure}

\begin{figure}[!t]
\centering
\includegraphics[width=\columnwidth,keepaspectratio]{figs/F6_relacional.png}
\caption{Relational analyses: (a) rate of effect-size reporting by research design; (b) digital-twin maturity (Kritzinger taxonomy) by publication year.}
\label{fig:F7}
\end{figure}

\subsection{Advanced relational analysis (exploratory)}\label{advanced-relational-analysis-exploratory}

Because per-theme frequencies cannot show how technical, pedagogical and
evaluative choices travel together, we examined pairwise associations between 18
binary study attributes spanning the four objectives (\ensuremath{\phi} coefficient with Fisher
exact tests over the 65 studies, plus non-parametric bootstrap 95\% confidence
intervals, B=5000).
With this sample size the associations are hypothesis-generating rather than
inferential; each is reported with its co-occurrence count and bootstrap CI, and
we flag edges whose interval includes zero. The attribute network (Figure 8) is
organised into three loosely connected communities rather than a single
integrated cluster: (i) a \textbf{technical-capability} group in which bidirectional
twins, real-time synchronisation and hybrid environments co-occur (\ensuremath{\phi}=0.69,
95\% CI {[}0.51, 0.86{]}, between bidirectional DT and real-time sync; \ensuremath{\phi}=0.40
{[}0.13, 0.65{]} between hybrid environment and bidirectional DT); (ii) an
\textbf{evaluation-rigour} group in which rigorous design, comparator use,
effect-size reporting and affective-outcome measurement co-occur (\ensuremath{\phi}=0.60
{[}0.42, 0.77{]} design--comparator; \ensuremath{\phi}=0.46 {[}0.22, 0.67{]} comparator--effect size);
and (iii) a \textbf{curricular} group linking undergraduate level, active learning,
industry collaboration and full life-cycle projects (\ensuremath{\phi}=0.43 {[}0.11, 0.70{]} between
life-cycle project and industry collaboration). One cross-objective signal is
worth noting: the technical-capability and curricular communities are bridged
mainly by a \emph{negative} association (active-learning-labelled courses are less
likely to use the most capable bidirectional twins). By contrast, the apparent
link between immersive laboratories and environmental-sustainability
quantification rests on only two studies and is not robust (\ensuremath{\phi}=0.29, 95\% CI
{[}\ensuremath{-}0.07, 0.60{]}, compatible with no association). Several attributes --- AI
augmentation, immersive environment, open-source licensing, sustainability
quantification --- remain isolated at the chosen threshold, indicating they do not
yet co-occur systematically with rigorous evaluation. Across the six highlighted
edges, five of six bootstrap CIs exclude zero, but several remain wide (e.g.
life-cycle--industry {[}0.11, 0.70{]}), reinforcing that the network is best read as
a map of hypotheses for confirmatory study.

The research-frontier gap map (Figure 9) plots each
empirical study on technical capability (a 0--1 composite of maturity,
synchronisation, AI and immersion) against evidence-quality score. Studies
concentrate in the low-capability columns (18 of 35) and in the lower
evidence-quality rows; the high-capability/strong-evidence corner --- studies that
are simultaneously technically advanced and rigorously evaluated --- contains only
1 study \cite{tarng_application_2024}. This near-empty frontier is the central structural gap of the field
and complements the qualitative decoupling reported above. The 3D evidence space
(Figure 10, with an interactive version) and the MCA embedding (Figure 11) tell
the same story from the multivariate side. The first MCA axis cleanly separates
empirical from non-empirical/conceptual studies, recovering a valid topological
structure of the field in which evaluation status is its dominant axis of
organisation. As is intrinsic to the geometry of MCA, the share of inertia carried
by this axis (12.3\%) is modest; this is a property of the method rather than a sign
of weak structure, and the separation should be read as a robust ordinal and
topological pattern rather than as a proportion of variance explained.

Finally, Tables 2 and 3 make the most important contributions explicit.
Table 2 lists the 15 empirical studies reaching at least a moderate
evidence-quality score, with their designs, samples, comparators and reported
effects --- including the controlled studies of \cite{kwateng_enhancing_2026}, between-group
gain +18.5 vs +4.3, t(198)=8.94, \cite{tarng_application_2024}, ANCOVA F=71.9, p\textless0.001,
\ensuremath{\eta^{2}}=0.50, \cite{lin_visfactory_2025}, +31.2\% transfer to novel problems with retention at
3 and 6 months, \cite{lu_digital_2025}, d=1.06 in algorithm design and
\cite{wang_innovative_2025}, +16.6\% vs +7.5\% knowledge mastery. These define the empirical frontier of
the field. Table 3 lists the four studies that assemble the most complete
technology--pedagogy--evidence stack \cite{sakkas_multiplayer_2025,tarng_application_2024,terkaj_framework_2024} or couple a detailed architecture with AI and evaluation
\cite{machado_automatic_2025}; notably, only \cite{tarng_application_2024} combines a high-maturity
twin with strong evidence, which is consistent with the decoupling and the empty
frontier identified above.

%\clearpage
%\onecolumn
\footnotesize
\begin{table}[!t]
	\renewcommand{\arraystretch}{1.2}
	% Reducción del espacio horizontal entre columnas nativo de LaTeX
	\setlength{\tabcolsep}{2.2pt} 
	\caption{Empirical Studies Reaching at Least Moderate Evidence Quality}
	\label{tab:T2}
	\centering
	\scriptsize
	\begin{tabular}{@{} l l l p{1.3cm} p{3.5cm} @{}}
		\hline
	
		\textbf{Study} & \textbf{Dom.} & \textbf{Design ($N$)} & \textbf{Comparator} & \textbf{Key Result, Score \& \{ES\}} \\ \hline
		\cite{kwateng_enhancing_2026} & Manuf. & Quasi-exp. (200) & Control (no int.) & \textbf{\{6\}} Pre/post +18.5 vs +4.3 (ctrl); $d=0.91$. \\
		\cite{lu_digital_2025} & Prog. & Quasi-exp. (135) & Trad. class & \textbf{\{6\}} Algorithm $d=1.06$; debug $d=0.91$. \\
		\cite{lin_visfactory_2025} & Autom. & Experim. (127) & Control group(s) & \textbf{\{5\}} +31.2\% transfer; higher retention at 3/6 m. \\
		\cite{tarng_application_2024} & Robot. & Quasi-exp. (75) & Trad. class & \textbf{\{5\}} Post 8.71 vs 3.75; ANCOVA $\eta^{2}=0.50$. \\
		\cite{wang_innovative_2025} & Cyber. & Quasi-exp. (60) & Trad. class & \textbf{\{5\}} Knowledge mastery +16.6\% vs +7.5\% (ctrl). \\
		\cite{khan_teaching_2026} & Manuf. & Mixed-m. (40) & Pre/post (W-S) & \textbf{\{4\}} 63\% improvement in evaluation scores. [ES: N] \\
		\cite{li_exploration_2025} & Chem. & Case/pilot (95) & Trad. class & \textbf{\{4\}} Knowledge gain 27.3\% vs 11.5\% (ctrl). \\
		\cite{liljaniemi_enhancing_2025} & Mech. & Quasi-exp. (96) & Physical lab & \textbf{\{4\}} VR group +11.85\% component ID. [ES: N] \\
		\cite{ning_exploration_2024} & Civil & Quasi-exp. (28) & Trad. class & \textbf{\{4\}} Significantly better performance. [ES: N] \\
		\cite{ortiz_digital_2025} & Autom. & Quasi-exp. (16) & Pre/post (W-S) & \textbf{\{4\}} Time 18.2$\rightarrow$12.5 min; error $\downarrow$. [ES: N] \\
		\cite{park_ar_2026} & NR & Quasi-exp. (10) & Trad. class & \textbf{\{4\}} Improved conceptual quiz scores. [ES: N] \\
		\cite{alvarez_ariza_investigating_2026} & Multi. & Case/pilot (119) & Pre/post (W-S) & \textbf{[4]} 96.5–100\% agreement on knowledge dev. \\
		\cite{garcia_digital_2026} & Logic & Case/pilot (30) & None & \textbf{\{3\}} Task completion 100\%/93\%/80\% by diff. \\
		\cite{speicher_digital_2026} & Manuf. & Case/pilot (6) & Trad. vs smart & \textbf{\{3\}} Machining time 58$\rightarrow$28 min ($>$50\% $\downarrow$). [ES: N] \\
		\cite{federal_university_of_amazonas_development_2025} & STEM & Mixed-m. (24) & Pre/post (W-S) & \textbf{\{3\}} Comprehension rose to 90.9\% (from 27\%). [ES: N] \\ \hline
		\multicolumn{5}{@{}l}{\parbox{8.8cm}{\vspace{1mm}\textit{Note:} Score in bold brackets \textbf{\{X\}} represents composite quality score out of 6. ES = Effect Size reported (assumed Yes unless explicitly tagged as \{ES: N\}). W-S = Within-Subject; NR = Not Reported.}}
	\end{tabular}
\end{table}


%\setlength{\tabcolsep}{4pt}
%\begin{longtable}{@{}
%  >{\raggedright\arraybackslash}p{(\linewidth - 14\tabcolsep) * \real{0.1100}}
%  >{\raggedright\arraybackslash}p{(\linewidth - 14\tabcolsep) * \real{0.1000}}
%  >{\raggedright\arraybackslash}p{(\linewidth - 14\tabcolsep) * \real{0.0850}}
%  >{\raggedright\arraybackslash}p{(\linewidth - 14\tabcolsep) * \real{0.0400}}
%  >{\raggedright\arraybackslash}p{(\linewidth - 14\tabcolsep) * \real{0.1150}}
%  >{\raggedright\arraybackslash}p{(\linewidth - 14\tabcolsep) * \real{0.0400}}
%  >{\raggedright\arraybackslash}p{(\linewidth - 14\tabcolsep) * \real{0.4200}}
%  >{\raggedright\arraybackslash}p{(\linewidth - 14\tabcolsep) * \real{0.0900}}@{}}
%\caption{Empirical studies reaching at least a moderate evidence-quality score (composite \ensuremath{\geq} 3 of 6), ordered by score. ES = effect size reported; NR = not reported. The composite sums design rigour, comparator, inferential statistics, effect size and declared validity threats.}\label{tab:T2}\\
%\toprule\noalign{}
%\begin{minipage}[b]{\linewidth}\raggedright
%Study
%\end{minipage} & \begin{minipage}[b]{\linewidth}\raggedright
%Domain
%\end{minipage} & \begin{minipage}[b]{\linewidth}\raggedright
%Design
%\end{minipage} & \begin{minipage}[b]{\linewidth}\raggedright
%N
%\end{minipage} & \begin{minipage}[b]{\linewidth}\raggedright
%Comparator
%\end{minipage} & \begin{minipage}[b]{\linewidth}\raggedright
%ES
%\end{minipage} & \begin{minipage}[b]{\linewidth}\raggedright
%Key reported result
%\end{minipage} & \begin{minipage}[b]{\linewidth}\raggedright
%Band (score)
%\end{minipage} \\
%\midrule\noalign{}
%\endhead
%\bottomrule\noalign{}
%\endlastfoot
%\cite{kwateng_enhancing_2026} & Manufacturing & Quasi-experiment & 200 & Control (no intervention) & Yes & Pre/post gain +18.5 vs +4.3 (control); t(198)=8.94, p\textless.001; d=0.91 (sensor-integration skills) & Strong (6) \\
%\cite{lu_digital_2025} & Programming & Quasi-experiment & 135 & Traditional class & Yes & Algorithm design d=1.06; debugging d=0.91; overall post-test d=0.82 & Strong (6) \\
%\cite{lin_visfactory_2025} & Automation & Experiment & 127 & Control group(s) & Yes & +31.2\% transfer to novel problems; higher retention at 3 and 6 months & Strong (5) \\
%\cite{tarng_application_2024} & Robotics & Quasi-experiment & 75 & Traditional class & Yes & Post-test 8.71 vs 3.75; ANCOVA F=71.9, p\textless.001, \ensuremath{\eta^{2}}=0.50 & Strong (5) \\
%\cite{wang_innovative_2025} & Cybersecurity & Quasi-experiment & 60 & Traditional class & Yes & Knowledge mastery +16.6\% vs +7.5\% (control) & Strong (5) \\
%\cite{khan_teaching_2026} & Manufacturing & Mixed-methods & 40 & Pre/post (within-subject) & No & 63\% improvement in course-evaluation scores & Moderate (4) \\
%\cite{li_exploration_2025} & Chemical eng. & Case study / pilot & 95 & Traditional class & Yes & Knowledge-mastery gain 27.3\% vs 11.5\% (control) & Moderate (4) \\
%\cite{liljaniemi_enhancing_2025} & Mechatronics & Quasi-experiment & 96 & Physical laboratory & No & VR group +11.85\% on component identification; parity on diagram interpretation & Moderate (4) \\
%\cite{ning_exploration_2024} & Civil / construction & Quasi-experiment & 28 & Traditional class & No & Significantly better academic performance & Moderate (4) \\
%\cite{ortiz_digital_2025} & Automation & Quasi-experiment & 16 & Pre/post (within-subject) & No & Task time 18.2\ensuremath{\rightarrow}12.5 min; error rate 21.6\%\ensuremath{\rightarrow}7.8\% & Moderate (4) \\
%\cite{park_ar_2026} & NR & Quasi-experiment & 10 & Traditional class & No & Improved conceptual-knowledge quiz scores (pilot) & Moderate (4) \\
%\cite{alvarez_ariza_investigating_2026} & Multi-case & Case study / pilot & 119 & Pre/post (within-subject) & Yes & 96.5--100\% agreement on fundamental-knowledge development & Moderate (4) \\
%\cite{garcia_digital_2026} & Digital logic & Case study / pilot & 30 & None & Yes & Task completion 100\% / 93.3\% / 80\% / 73.3\% by difficulty & Moderate (3) \\
%\cite{speicher_digital_2026} & Manufacturing & Case study / pilot & 6 & Traditional vs smart-mfg & No & Machining cycle time 58\ensuremath{\rightarrow}28 min (\textgreater50\% reduction) & Moderate (3) \\
%\cite{federal_university_of_amazonas_development_2025} & Postgraduate STEM & Mixed-methods & 24 & Pre/post (within-subject) & No & Share rating comprehension 4--5 rose to 90.9\% (from 27.3\%) & Moderate (3) \\
%\end{longtable}
\normalsize
%\twocolumn


%\clearpage
%\onecolumn
%\footnotesize
%\setlength{\tabcolsep}{4pt}
%\begin{longtable}{@{}
%  >{\raggedright\arraybackslash}p{(\linewidth - 14\tabcolsep) * \real{0.1200}}
%  >{\raggedright\arraybackslash}p{(\linewidth - 14\tabcolsep) * \real{0.1500}}
%  >{\raggedright\arraybackslash}p{(\linewidth - 14\tabcolsep) * \real{0.1000}}
%  >{\raggedright\arraybackslash}p{(\linewidth - 14\tabcolsep) * \real{0.0750}}
%  >{\raggedright\arraybackslash}p{(\linewidth - 14\tabcolsep) * \real{0.1650}}
%  >{\raggedright\arraybackslash}p{(\linewidth - 14\tabcolsep) * \real{0.0900}}
%  >{\raggedright\arraybackslash}p{(\linewidth - 14\tabcolsep) * \real{0.2000}}
%  >{\raggedright\arraybackslash}p{(\linewidth - 14\tabcolsep) * \real{0.1000}}@{}}
%\caption{Reference-architecture configurations: studies assembling the most complete technology--pedagogy--evidence stack (full-stack T1+T2+T5+T6) or coupling a detailed architecture with AI and empirical evaluation (T2+T3+T6). DT type follows the Kritzinger et al.~taxonomy; NR = not reported.}\label{tab:T3}\\
%\toprule\noalign{}
%\begin{minipage}[b]{\linewidth}\raggedright
%Study
%\end{minipage} & \begin{minipage}[b]{\linewidth}\raggedright
%Configuration
%\end{minipage} & \begin{minipage}[b]{\linewidth}\raggedright
%Environment
%\end{minipage} & \begin{minipage}[b]{\linewidth}\raggedright
%DT type
%\end{minipage} & \begin{minipage}[b]{\linewidth}\raggedright
%AI
%\end{minipage} & \begin{minipage}[b]{\linewidth}\raggedright
%Pedagogy
%\end{minipage} & \begin{minipage}[b]{\linewidth}\raggedright
%System modelled
%\end{minipage} & \begin{minipage}[b]{\linewidth}\raggedright
%Evidence band
%\end{minipage} \\
%\midrule\noalign{}
%\endhead
%\bottomrule\noalign{}
%\endlastfoot
%\cite{sakkas_multiplayer_2025} & Full-stack (T1+T2+T5+T6) & Virtual-only & Digital shadow & --- & Active learning & Electronic circuits (Arduino, ESP32) & Weak \\
%\cite{tarng_application_2024} & Full-stack (T1+T2+T5+T6) & Hybrid physical--virtual & Digital twin & --- & NR & WLKATA six-axis industrial robotic arm & Strong \\
%\cite{terkaj_framework_2024} & Full-stack (T1+T2+T5+T6) & Hybrid physical--virtual & NR & --- & NR & Manufacturing plant & Weak \\
%\cite{machado_automatic_2025} & Architecture + AI + evidence (T2+T3+T6) & Remote & NR & CNN, generative AI, LLM, data-driven & NR & Manufacturing plant (CP-LAB) & Weak \\
%\end{longtable}
\begin{table}[!t]
	\renewcommand{\arraystretch}{1.2}
	% Ajuste micro-espacial de columnas nativo de la tabla
	\setlength{\tabcolsep}{2.5pt} 
	\caption{Reference-Architecture Configurations and Full-Stack Couplings}
	\label{tab:T3}
	\centering
	\scriptsize
	\begin{tabular}{@{} l l l p{3.6cm} c @{}}
		\hline
		\textbf{Study} & \textbf{Config.} & \textbf{Platform / Type} & \textbf{System Modelled \& AI Tools} & \textbf{Evid.} \\ \hline
		\cite{sakkas_multiplayer_2025} & T1+T2+T5+T6 & Virtual / Shadow & Electr. circuits (Arduino, ESP32)\textmd{\dag}. \textit{AI:} None. & Weak \\
		\cite{tarng_application_2024} & T1+T2+T5+T6 & Hybrid / Twin & WLKATA 6-axis industrial robot arm. \textit{AI:} None. & Strong \\
		\cite{terkaj_framework_2024} & T1+T2+T5+T6 & Hybrid / NR & Manufacturing plant. \textit{AI:} None. & Weak \\
		\cite{machado_automatic_2025} & T2+T3+T6 & Remote / NR & Manufacturing plant (CP-LAB). \textit{AI:} CNN, GenAI, LLM. & Weak \\ \hline 
		\multicolumn{5}{@{}l}{\parbox{8.8cm}{\vspace{1mm}\textit{Note:} DT type follows Kritzinger et al. taxonomy; NR = Not Reported. Stack components: T1 = practical infrastructure (environment), T2 = enabling architecture, T3 = AI convergence, T5 = pedagogical design, T6 = empirical evidence.}}\\
		\multicolumn{5}{@{}l}{\textmd{\dag} Utilizes active learning pedagogy.}
	\end{tabular}
\end{table}
\normalsize
%\twocolumn


\begin{figure}[!t]
\centering
\includegraphics[width=\columnwidth,keepaspectratio]{figs/F7_red_atributos.png}
\caption{Attribute-association network. Nodes are binary study attributes; edges are pairwise \ensuremath{\phi} associations whose bootstrap 95\% confidence interval (CI) excludes zero (B = 5,000). Communities: technical capability, evaluation rigour, curricular design.}
\label{fig:F8}
\end{figure}

\begin{figure}[!t]
\centering
\includegraphics[width=\columnwidth,keepaspectratio]{figs/F8_mapa_brechas.png}
\caption{Capability-by-evidence gap map: each empirical study plotted by technical-capability composite (0--1: maturity, synchronisation, AI, immersion) against evidence-quality score (0--6). The high-capability, strong-evidence corner is near-empty.}
\label{fig:F9}
\end{figure}

\begin{figure}[!t]
\centering
\includegraphics[width=\columnwidth,keepaspectratio]{figs/F9_espacio_3d.png}
\caption{Three-dimensional evidence space (technical capability \ensuremath{\times} evidence quality \ensuremath{\times} publication year). An interactive version is provided as supplementary material.}
\label{fig:F10}
\end{figure}

\begin{figure}[!t]
\centering
\includegraphics[width=\columnwidth,keepaspectratio]{figs/F10_mca_embedding.png}
\caption{Multiple correspondence analysis (MCA) embedding of study attributes. The first axis (12.3\% of inertia) separates empirical from non-empirical/conceptual studies.}
\label{fig:F11}
\end{figure}

\subsection{Summary}\label{summary}

The corpus consistently demonstrates feasibility and engagement and
is disciplinarily broad, with a recent shift toward more capable, bidirectional,
and occasionally AI-augmented twins. The weight of the evidence, however, rests
on a minority of methodologically strong studies; competency and benefit claims
frequently lack matched empirical or quantitative support; and the technical
description of the twin is largely separate from the evaluation of its
educational effect. These patterns directly inform the gaps and future-research
agenda developed in the Discussion.

\section{Discussion}\label{discussion}

This review set out to characterise educational digital twins (DTs) across four
objectives and, crucially, to examine how their technical, pedagogical and
evaluative dimensions travel together. The synthesis of 65 studies yields a
consistent but sobering picture: the field reliably demonstrates feasibility and
engagement, is disciplinarily broad, and is shifting toward more capable,
bidirectional and occasionally AI-augmented twins, yet its evidence rests on a
small minority of methodologically strong studies. Interpreting this pattern, not
the headline enthusiasm, is the purpose here.

The most consequential finding is structural rather than aggregate: across the
corpus, the maturity of the twin is largely decoupled from the rigour with which
it is evaluated. The studies that describe how complete a twin is---its
bidirectional synchronisation, its layered architecture---are seldom the studies
that test whether it improves learning, and the capability-by-evidence frontier
is almost empty. Read through constructive alignment \cite{biggs_enhancing_1996,biggs_teaching_2011}, this is a systemic misalignment: the field has invested
heavily in one element of the instructional system---the learning activity, the
twin itself---while the intended learning outcomes and the assessment that would
evidence them remain comparatively underspecified. The pattern is, in this sense,
a familiar one in engineering and STEM education, where the adoption of a new
instructional technology routinely outpaces the rigorous evaluation of its effect
\cite{henderson_facilitating_2011}, and it reproduces for digital twins the very
situation that comparative reviews of virtual and remote laboratories have
documented for two decades: enthusiastic deployment, a prevailing ``no significant
difference'' verdict, and a thin and methodologically weak evidence base beneath it
\cite{brinson_learning_2015,heradio_virtual_2016}. This decoupling also mirrors a
tension documented across the wider DT
literature, where academic sophistication consistently outpaces validated
deployment and workforce readiness \cite{chen_ai-enhanced_2025}. It is compounded by
terminological inflation: the prevalence of self-described ``complete'' twins,
rarely accompanied by latency or synchronisation measurements, suggests that the
label is applied more loosely than the established physical-to-virtual and
virtual-to-physical data-flow criteria would warrant \cite{thelen_comprehensive_2022}, and
that the field still lacks shared criteria for evaluating what a twin actually
does \cite{liu_cognitive_2025}. Claims of Industry 4.0/5.0 competence and institutional
benefit similarly outpace demonstrated outcomes, and the asymmetry between
benefit-only and barrier-only studies counsels caution in reading the field's
optimism at face value.

One alternative reading of the decoupling deserves stating, because it tempers a
purely critical interpretation. The separation between capable-but-unevaluated and
evaluated-but-simple twins partly tracks a division of publication labour: in our
corpus, architecture-heavy work appears disproportionately in conference and
engineering venues, whereas the studies that mount controlled evaluations tend to
be journal articles, and conference papers score markedly lower on the
evidence-quality composite than journal articles. Part of what looks like a
within-study failure to evaluate may therefore reflect where each kind of
contribution is published. This does not dissolve the problem---the field as a whole
still lacks studies that are both technically complete and rigorously
evaluated---but it locates it at the level of the field's reporting ecosystem rather
than of individual authors' competence.

Three exploratory analyses sharpen this reading. The thematic co-occurrence
structure (Figure 6) shows a field organised around a pedagogy--competencies--
evidence core---competencies and pedagogy co-occur in 34 studies and pedagogy and
evidence in 31---while the AI dimension stays peripheral, never co-occurring with
another theme in more than eight studies; the DT--AI convergence prominent in the
industrial literature \cite{chen_ai-enhanced_2025,liu_cognitive_2025} thus remains marginal
and largely unvalidated in education. The
attribute-association network (Figure 8) resolves the studies into three loosely
connected communities---technical capability, evaluation rigour and curricular
design---bridged between the technical and curricular groups mainly by a \emph{negative}
association, whereby courses labelled ``active learning'' tend not to employ the
most capable bidirectional twins. The research-frontier gap map (Figure 9)
localises the decoupling: of 35 empirical studies, eighteen fall in the lowest-capability
columns and only one occupies the high-capability, strong-evidence cell, while the
multivariate embedding (Figure 11) recovers a valid topological structure whose
dominant axis separates empirical from non-empirical work. Tellingly, effect-size
reporting is flat across design types even as twin maturity rises over time
(Figure 7), so weak quantification is a field-wide habit, not an artefact
of the weakest designs.

These patterns demand interpretive caution. They derive from 65 studies spread
across many sparsely populated cells; the association estimates carry wide
bootstrap confidence intervals, several compatible with no effect---the apparent
link between immersive laboratories and sustainability quantification, for
instance, rests on only two studies---and the embedding accounts for only a modest
share of inertia. They are therefore best read as a map of hypotheses for
confirmatory study, not as established relationships. That such an exploratory
layer is necessary is itself diagnostic of how little well-reported empirical work
the field has produced.

These observations extend, rather than echo, prior syntheses of DTs in education.
Earlier bibliometric work mapped the field's rapid growth and flagged unresolved
questions around teacher competence and governance, but did not weigh the quality
of the underlying learning evidence \cite{chamorro-atalaya_use_2024};
domain-specific reviews reached a ``strong consensus'' on DT effectiveness for
training while themselves noting the fragmented, largely conceptual nature of the
literature \cite{omrany_digital_2026}; and materials-science mapping documents
exponential output without appraising educational outcomes at all
\cite{vergara_decade_2025}. By coupling a structured description of architecture and
pedagogy with an explicit, sensitivity-tested quality appraisal and a relational
analysis, the present review adds the dimension these accounts lack: not how much
is published, but how well what is claimed has been shown.

The few strong studies in our corpus report substantial gains, of the order of
Cohen's d \ensuremath{\approx} 0.9--1.25. These sit at, or above, the upper bound of
meta-analytic estimates for comparable immersive interventions---knowledge-
acquisition effects around a standardised mean difference of 0.33--0.48 in
virtual- and augmented-reality simulation \cite{li_effects_2026,almekkawi_effectiveness_2026}, and collaborative- or mobile-VR effects near Hedges'
g 0.47--0.50 \cite{chen_effects_2026,widowati_meta-analysis_2025}. This positioning warrants
restraint. Such meta-analyses report high heterogeneity and inconsistent effects
across domains; more pointedly, effect sizes in this literature have declined
over time as designs have become
more rigorous---from very large early estimates to small recent ones
\cite{widowati_meta-analysis_2025}---a trajectory consistent with novelty effects and weak
controls. Our predominantly recent, single-institution, non-randomised studies
are precisely those most exposed to this inflation, reinforcing that the reported
gains are promising but not yet robustly generalisable.

A second under-examined frontier concerns whether learning on a twin persists and
transfers to physical practice---the upper, more consequential levels of
Kirkpatrick's evaluation hierarchy \cite{kirkpatrick_kirkpatricks_2016}. The corpus
sits overwhelmingly at the lower levels: reactions and immediate post-tests
dominate, whereas to our reading only one study measured learning retention at a
specified multi-month follow-up (three and six months) \cite{lin_visfactory_2025}, and
none measured transfer to subsequent physical-laboratory or workplace performance.
The broader evidence offers little reassurance: in a review of 64 VR studies, only
two examined transfer at all \cite{zhong_virtual_2026}. Transfer cannot be assumed from immersion or visual
fidelity; it depends on functional rather than merely perceptual correspondence
\cite{harris_framework_2020}, and inadequate technical fidelity can even produce
negative transfer \cite{kim_effect_2026}. The corpus's habit of asserting twin
``completeness'' without synchronisation or fidelity metrics is therefore not a
minor reporting gap but a direct threat to the ecological-validity claims on which
much of the pedagogical case rests.

Several implications follow. For curriculum design, DTs should be paired with an
intentional, named pedagogy and aligned assessment rather than deployed as
active-learning settings whose instructional model is left implicit. Concretely,
future implementations should move beyond generic ``active experimentation'' and
anchor their designs in explicit learning theory: constructivist and
social-constructivist principles for knowledge building, Biggs' constructive
alignment to lock intended competencies, twin-based activities and assessment into
a single coherent system \cite{biggs_enhancing_1996}, and---where the goal is learner
self-direction in open-ended, capability-oriented tasks---student-agency models such
as heutagogy \cite{blaschke_heutagogy_2012}. One design in particular is conspicuously
under-exploited: having students \emph{construct} the twin rather than merely operate a
ready-made one. In ICAP terms this is the apex of cognitive engagement
\cite{chi_icap_2014}, and in constructionist terms it makes the learner the builder of
the model they are meant to understand; yet our corpus documents it only once,
even though the learning-factory tradition it most naturally belongs to is already
present in the field. The negative technical--curricular bridge
in the attribute network suggests engineering ambition and pedagogical design are
currently advancing on parallel tracks; naming and theorising the pedagogy is what
would allow the instructional design to keep pace with the sophistication of the
twin. For
researchers, the priority is methodological, and it is the one the
engineering-education research community has urged for two decades: holding studies
of instructional innovation to the standards of rigorous educational research
\cite{streveler_conducting_2006}. Concretely, this means
pre-registered comparators, reported effect sizes, and explicit measurement of
retention and transfer, supported by designs---including rigorous single-case and
quasi-experimental approaches---suited to authentic laboratory settings
\cite{dayo_evaluating_2024}, with the reporting transparency that educational-technology
research is repeatedly found to lack \cite{demir_analyzing_2025}. For institutions,
accessibility and sustainability claims need instrumentation rather than
assertion, particularly where digital-divide constraints shape who benefits
\cite{eltaiba_benefits_2025}. For tool builders,
interoperability, synchronisation metrics and the validation of AI components
remain prerequisites for credible educational deployment \cite{chen_ai-enhanced_2025}.

Several limitations qualify this synthesis. The corpus is recent and skewed
toward manufacturing and automation and undergraduate cohorts; the
relational analyses rest on small cell counts and are hypothesis-generating rather
than confirmatory; and the evidence-quality composite is necessarily ad hoc,
though its banding proved robust under threshold-sensitivity analysis. Full-text
attrition was substantial: 61 of 126 eligible records could not be retrieved, and
this loss was differential by document type---82\% of the unretrieved reports were
conference papers, against 40\% of the included corpus. Because conference papers in
our corpus score markedly lower on evidence quality than journal articles, the
missing reports would, if anything, be expected to reinforce rather than soften the
weak-dominated profile we report, so the corpus median is best read as a
conservative bound; nonetheless the gap limits coverage of conference-reported
architectures, and of 2024 in particular. Relatedly, we did not formally assess
publication or reporting bias---no meta-analytic funnel test is meaningful here---but
the asymmetry between benefit-only and barrier-only studies is consistent with
selective reporting and should be read in that light. Finally, coverage of 2026 is
necessarily partial, since the search was run mid-year and early-access records
continue to appear. We
deliberately refrained from meta-analysis, a decision the marked clinical and
methodological heterogeneity of the broader literature retrospectively supports
\cite{widowati_meta-analysis_2025}. The agenda these gaps define is concrete: populate the
empty capability-by-evidence frontier with pre-registered, multi-institution
quasi-experiments on high-maturity twins; measure transfer and retention rather
than immediate satisfaction; validate AI-augmented and cognitive twins against
declared performance criteria; and standardise the reporting of synchronisation,
cost and sustainability, so that the field's momentum can mature into cumulative,
trustworthy evidence.

\section{Conclusions}\label{conclusions}

Regarding the technological architectures and deployment environments (O1),
educational digital twins are realised predominantly as hybrid physical--virtual
and immersive configurations rather than as fully instrumented, closed-loop
systems. The field's technical ambition is rising---bidirectional and occasionally
AI-augmented twins are increasingly described---but the terminology frequently
outpaces verified synchronisation, and the convergence with artificial
intelligence remains incipient and largely unvalidated. The technological maturity
of the field is therefore, at present, more often asserted than demonstrated, and
claims of completeness are best read as descriptive rather than evidential until
synchronisation and fidelity are actually measured.

On the pedagogical and curricular dimension (O2), digital twins are consistently
embedded in active, experiential and project-based designs and span a widening
range of engineering disciplines, signalling pedagogical intent and
disciplinary breadth. Yet the instructional model is often left implicit, and the
competencies attributed to Industry 4.0/5.0 are seldom anchored to recognised
standards or validated with industry. The pedagogical case is thus plausible but
under-specified: technical capability and intentional instructional design are
advancing on parallel rather than integrated tracks.

With respect to the empirical evidence of learning outcomes (O3), the corpus
establishes feasibility and positive engagement, and a minority of methodologically
stronger studies report meaningful gains; these, however, sit within a
heterogeneous and weakly reported base in which inferential statistics and effect
sizes are the exception rather than the norm. The present evidence justifies
cautious optimism about engagement and proof of concept, but does not yet support
robust or generalisable claims of learning efficacy, particularly concerning skill
transfer and longer-term retention.

Concerning institutional benefits and barriers (O4), accessibility, educational
continuity and sustainability are argued far more often than they are quantified,
while the technical, pedagogical and institutional barriers are reported
unevenly. The presence of a substantial barriers-aware literature indicates that
the field is not uniformly promotional; even so, the recurring assertion of
benefits without a matched account of limitations counsels restraint. The
institutional value proposition remains largely rhetorical pending the
instrumentation of cost, environmental and equity claims.

Taken together (general objective), this review characterises a field that
reliably demonstrates feasibility, engagement and disciplinary breadth, yet whose
defining feature is a structural decoupling between the maturity of the twin and
the rigour of its evaluation, with competence and institutional claims outpacing
demonstrated outcomes. The contribution of this synthesis is not to adjudicate
whether digital twins ``work'', but to weight the available evidence and make this
decoupling---and the near-empty frontier where technical sophistication meets strong
evaluation---explicit. Digital twins in engineering education are best understood as
a promising but not yet evidentially consolidated approach, whose trajectory will
depend less on further technical demonstration than on the quality of the evidence
brought to bear on it. In this, the review offers engineering education research a
sharper criterion for judging progress---evidential quality rather than technical
novelty---and offers practice a caution against equating sophistication with
demonstrated learning. Read against the longer arc of laboratory-education
research, the digital twin has so far inherited the comparative question that
hands-on, virtual and remote laboratories have faced since \cite{ma_handson_2006} and
\cite{brinson_learning_2015}---whether a technology-mediated laboratory achieves the
learning outcomes of physical practice---without yet supplying the rigorous evidence
needed to answer it. The broader lesson the field can carry forward is therefore
methodological as much as technological: the evaluation of a new instructional
technology must be designed into it from the outset, as an aligned system of
intended outcomes, activities and assessment, rather than appended once the
artefact is built.

Several recommendations follow, addressed first to the communities that shape what
counts as evidence. For educators, digital twins should be paired with an explicit
pedagogical model and aligned assessment rather than deployed as implicitly
``active'' settings: the intended competencies, the twin-based activities and the
assessment should be designed as a single aligned system \cite{biggs_teaching_2011},
the activities pitched at the constructive and interactive levels of cognitive
engagement that predict the largest gains \cite{chi_icap_2014}, and the targeted
competencies anchored to the recognised objectives of the instructional
laboratory \cite{feisel_role_2005} and to accreditation and curricular frameworks such
as ABET and CDIO rather than asserted as generic ``Industry 4.0/5.0'' readiness. For
researchers, the priority is evidential consolidation: pre-registered,
multi-institution studies on higher-maturity twins, with explicit comparators---
including, where the claim is to substitute for hands-on work, a physical-laboratory
comparison---reported effect sizes, and the measurement of transfer and retention
rather than immediate satisfaction; AI and synchronisation components should be
validated against declared performance criteria. For institutions and, more
narrowly, tool developers, cost, sustainability, accessibility, interoperability
and synchronisation should be reported with metrics rather than asserted. To
support cumulative, verifiable progress, the full protocol, screening trace,
extraction dataset and analysis code underlying this review are released openly
(GitHub and Zenodo), so that the synthesis can be audited, reused and extended.
Before any update, the review itself would be strengthened by registering the
protocol and recovering the full texts that could not be retrieved, so that the
corpus on which these conclusions rest is as complete as the field allows.
%
%\section{Statement on Artificial Intelligence}\label{statement-on-artificial-intelligence}
%
%In the interest of transparency, and consistent with the journal's policy, we
%disclose the following uses of artificial intelligence (AI) in this work, beyond
%the standard copyediting tools embedded in word processors (e.g., grammar
%checkers).
%
%\emph{Inductive thematic analysis.} The eight inductive themes (T1--T8) were derived
%with three frontier LLMs --- ChatGPT 5.5, Gemini 3.1 Pro, and Claude Opus 4.7 ---
%whose independent codings were iteratively cross-fed and refined under author
%review. The resulting theme--objective scheme and all decisions were determined by
%the authors; the three models' raw outputs are available in the public repository.
%
%\emph{Data screening.} Three independent large language model (LLM) agents --- implemented
%with DeepSeek-V4-Flash in a non-reasoning configuration at temperature 0 to ensure
%deterministic, reproducible classification --- performed the semantic screening of
%records against the protocol's pre-specified inclusion/exclusion tests and
%construct-severity mapping, each returning a decision, a confidence score, and the
%governing criterion; a fourth agent (same model and settings) acted as moderator
%for split or uncertain cases. AI-assigned decisions were not treated as final:
%inter-agent reliability was quantified (Fleiss' \ensuremath{\kappa} = 0.60), all records flagged as
%include or uncertain were reviewed by the authors, and the eligibility of every
%included study was confirmed by a human. This use is described in the Methods
%(Study selection).
%
%\emph{Data processing.} Full-text PDFs were converted to structured text with an
%automated document-parsing pipeline (Marker) to support extraction. Every
%extracted value was verified by a human against the source and recorded with a
%verbatim supporting quotation and its location (Methods, Data extraction). The
%analytical results --- descriptive profiling, the evidence-quality composite, and
%the relational analyses --- were produced by an author-developed, reproducible
%computational pipeline, not generated by AI.
%
%\emph{Manuscript development.} Generative AI tools were used to assist with drafting
%and language refinement of portions of the manuscript. All AI-assisted text was
%reviewed, edited, and verified by the authors.
%
%AI tools were not used to generate, alter, or fabricate data; to select,
%invent, or format references; or to produce the study's findings. The authors
%take full responsibility for the validity, accuracy, and integrity of the entire
%manuscript, including all claims, citations, and references.
%
%\section{Data and Code Availability}\label{data-and-code-availability}
%
%All materials supporting this review are openly available: the review protocol and
%extraction schemas, the database search exports, the full screening trace
%(per-record agent decisions, confidence scores, governing criteria, and final
%status), the extraction dataset with verbatim supporting quotations, the
%evidence-quality scoring, and the complete, end-to-end analysis and
%figure-generation code --- including the theme co-occurrence, multiple correspondence
%analysis, and attribute-association-network pipelines. These are released in a
%public GitHub repository and archived on Zenodo with a citable DOI ({[}repository URL
%and Zenodo DOI to be inserted on acceptance{]}), under an MIT license for code and a
%CC BY 4.0 license for data and text. The review followed a pre-specified protocol
%(Review Protocol v2); the protocol will be registered (PROSPERO/OSF) before any
%update.


\bibliographystyle{IEEEtran}
\bibliography{master}

\end{document}
